\documentclass{article}
\usepackage{amsmath}
\usepackage{amssymb}
\usepackage{graphicx}
\usepackage[options]{mcode}

\newcommand{\tr}{^\mathsf{T}}

\author{Jeff Steward}
\title{Introduction to Particle Filtering}
\begin{document}
\maketitle

\section{Introduction to filters}

A ``filter'' in this context is used to describe a mathematical process 
that removes an unwanted component from data.  This could be thought of as 
removing noise from data, removing some frequencies of the data (through a 
Fourier-transform, for example), or smoothing the data.

For our purposes, we will take a filter to be software that attempts 
to remove noise from data.  We can look at this either from the perspective
of stochastic differential equations or from a probability perspective.  In 
this tutorial, I will attempt to show both approaches as they are equally 
useful.

\subsection{An example}

We will consider the following situation.  First, suppose we have a model of how
a system evolves that includes a random component.  For example, we could assume
that there is a car that is driving around in a parking lot with the driver 
constantly choosing random acceleration.  The model for this could look like:

\[ \begin{array}{rcl}
dx/dt & = & u \\
dy/dt & = & v \\
du/dt & = & a_x \\
dv/dt & = & a_y
\end{array}\] 

Where $(x,y)$ is the position and $(u,v)$ are the $x$ and $y$ component of
acceleration.  We could assume that $(a_x,a_y)$ is ``random'' - that is, it
is changing all the time.  

This description, while intuitive, is unfortunately insufficent mathematically.
What is meant by ``changing all the time''?  We therefore use the more 
complicated-looking stochastic differential equation model

\[\begin{array}{rcl}
dx_t & = & u_t~dt \\
dy_t & = & v_t~dt \\
du_t & = & a_x~dB_t \\
dv_t & = & a_y~dB_t
\end{array}\]

Here, $\mathbf{x}_t = (x_t,y_t,u_t,v_t)\tr$ is a \emph{stochastic process}, 
meaning that at each time $t_k$, $x_k,y_k,u_k$ and $v_k$ are random variables. 
While a random variable is a function that maps an event to a probability,
a stochastic process is a function that maps a time index to a random variable. 
The $B$ here stands for Brownian motion, and $dB_t$ means that the variance of 
the random variables $u_t$ and $v_t$ grows by an amount $a_x^2~dt$ and 
$a_y^2~dt$ in time $dt$.

To say this in a more understandable way, the SDE

\[dx_t = \mu(x_t,t)dt + \sigma(x_t,t)dB_t\]
means that in a small amount of time $\delta t$, 
the random variable $x_{t + \delta t}$ will have changed its mean, on average,
by $\mu(x_t,t) \delta t$ when compared to $x_t$.  To be precise, the change in 
mean of $x_{t + \delta t}$ will be normally distributed with mean 
$\mu(x_t,t)\delta t$ and variance $\sigma^2(x_t,t)$.  

This interpretation leads us to the discrete version of the system above.  
We discretize the equation at $N_t$ regularly spaced points in time
$t = (0,t_f)$, and we call $(x_k^d,y_k^d,u_k^d,v_k^d)$ our 
discrete values
that approximate the value of $(x_k,y_k,u_k,v_k)$ at time $t_k = k \Delta t$, 
where $\Delta t = t_f/N_t$.  From here on, we will drop the superscript $d$ on 
the discrete variables as it is tedious to write and we are most interested in 
the discrete case.

If we approximate $dx_k = x_k - x_{k-1}$, we find

\[\begin{array}{rcl}
x_k & = & x_{k-1} + u_k \Delta t \\
y_k & = & y_{k-1} + v_k \Delta t \\
u_k & = & u_{k-1} + \eta_{x_k} \Delta t \\
v_k & = & v_{k-1} + \eta_{y_k} \Delta t
\end{array}\]  

Here, 
$\eta_{x_k}$ and $\eta_{y_k}$ are normally distributed variables with 
mean 0 and variances of $a_x^2$ and $a_y^2$, respectively.  Since $x_k$ and 
$y_k$ involve $u_k$ and $v_k$ from the same time step, we substitute in the 
value for $u_k$ and $v_k$ in those expressions, giving

\[\begin{array}{rcl}
x_k & = & x_{k-1} + u_{k-1} \Delta t + \eta_{x_k} \Delta t^2\\
y_k & = & y_{k-1} + v_{k-1} \Delta t + \eta_{y_k} \Delta t^2\\
u_k & = & u_{k-1} + \eta_{x_k} \Delta t \\
v_k & = & v_{k-1} + \eta_{y_k} \Delta t
\end{array}\]  

We now state our system in terms of probability.  We recall that for a normal
distribution $X \sim \mathcal{N}(\mu, \sigma^2)$, the variable 
$Y = aX + b \sim \mathcal{N}(a\mu + b, a^2\sigma^2)$.  Since the noise term is 
normally distributed, we see that

\[
\begin{array}{rcl}
x_k & \sim & \mathcal{N}(x_{k-1} + u_{k-1} \Delta t, a_{x_k}^2 \Delta t^4) \\
y_k & \sim & \mathcal{N}(y_{k-1} + v_{k-1} \Delta t, a_{y_k}^2 \Delta t^4) \\
u_k & \sim & \mathcal{N}(u_{k-1}, a_{x_k}^2 \Delta t^2) \\
v_k & \sim & \mathcal{N}(v_{k-1}, a_{y_k}^2 \Delta t^2)
\end{array}\]

Since the noise terms $\eta_x$ and $\eta_y$ are now present, we cannot 
hope to ``solve'' the system in the sense of finding ``the'' solution as we 
can for an ordinary differential equation or a partial 
differential equation.  Instead, we will seek the probability density of the 
possible states at each step in time.  In order to
find this probability density function, we could analyze the problem directly
if the system is relatively simple, or we could discretize the problem and run 
a numerical version of the system several times to learn from the statistics of
our runs.  

When we run the system many times, we call a collection of runs of a system 
an ``ensemble'' and each individual run a ``member'' or, in the case of 
particle filters, a ``particle.''  From the ensemble, we will be able to
retrieve emperical statistics that will give us the approximate probability 
density function of states.

\subsection{Observations}

We now introduce the more realistic case where instead of knowing the state 
of the system directly, we have observations of the system that have noise
in them.  For example, imagine that in the car problem considered above you 
have some measurements of the $x$ and $y$ location of the car at each time
step $k$.  The machine you are getting these observations from is not 
perfect and the observations have some error in them.  For the sake of 
simplicity, we wish to consider the case where the errors are normally 
distributed.  Let's also assume we know the mean and variance of these errors.
In reality, one would of course hope the mean was zero (i.e. the measurements 
are unbiased) and the variance of the error is small.

We want to take these observations and estimate what the state of the system
should be.  We will use our model of the system (which we will trust is an 
accurate depiction of the system) and the known statistics of the state change
and observation error in order to make a prediction of the state $\mathbf{x}_k$
at each time step $k = 1,2,\ldots,N_t$.  This prediction will depend on the 
quality of the observations we have, i.e. on the variance of the observation 
errors.

We create a model of the observations as a function of the state, which 
in this case is very simple.

\[\begin{array}{rcl}
z_{k_x} & = & x_k + n_{k_x} \\
z_{k_y} & = & y_k + n_{k_y}
\end{array}\]

Here, $\mathbf{z}_k = (z_{k_x},z_{k_y})\tr$ are our imperfect observations.  
$x_k$ and $y_k$ represent the true state, which it is assumed our state model 
accurately represents.  $n_{k_x}$ and $n_{k_y}$ are Gaussian variables with
mean $0$ and variance $R_k$. 

Our task is to estimate $\mathbf{x}_k$ based on the observations $\mathbf{z}_k$,
our model of the system, knowing the variance of the state change $Q_k$, 
our model of the observations, and the variance of our observations $R_k$.  We
want to do this at each time step $k$, although we assume we only have the 
information for time steps $i \le k$.  This is a very important problem in 
many real world situations.

\subsection{Generalization}

We now generalize the problem described above.  Suppose that we have a state
model

\[\mathbf{x}_k = f_k(\mathbf{x}_{k-1}, \mathbf{v}_{k-1})\]

where $\mathbf{x}_k, \mathbf{x}_{k-1} \in \mathbb{R}^{N_x}$ are the state
variables at time $k$ and $k-1$, respectively, 
$\mathbf{v}_{k-1} \in \mathbb{R}^{N_{v}}$ is 
multivariate Gaussian noise with mean $\{\mathbf{0}\}_{\mathbb{R}^{N_v}}$ and 
covariance $Q_k \in \mathbb{R}^{N_v} \times \mathbb{R}^{N_v}$.  Note that the 
covariance of the noise can change with $k$.  $f_k$ can be a non-linear 
function, and can also change at each time step $k$.

Suppose we also have an observation model

\[\mathbf{z}_k = h_k(\mathbf{x}_{k}, \mathbf{n}_k)\]
where $\mathbf{z}_k \in \mathbb{R}^{N_z}$ is the observation at time step $k$, 
$\mathbf{x}_k \in \mathbb{R}^{N_x}$ is the state at time $k$, and 
$\mathbf{n}_k$ is multivariate Gaussian noise with mean 
$\{\mathbf{0}\}_{\mathbb{R}^{N_n}}$ and covariance 
$R_k \in \mathbb{R}^{N_n} \times \mathbb{R}^{N_n}$.  $h_k$ is a possibly
non-linear function.  Both $R_k$ and $h_k$ are allowed to vary at each time
step $k$.

In this framework, our example problem becomes

\[\mathbf{x}_k = f_k(\mathbf{x}_{k-1}, \mathbf{v}_{k-1}) = 
F\mathbf{x}_{k-1} + G\mathbf{v}_{k-1}\]

where $\mathbf{x}_k = (x,y,u,v)\tr$, 
$\mathbf{v}_{k-1} = (\eta_x,\eta_y)\tr$ is a 
multivariate Gaussian variable with mean $\{\mathbf{0}\}_{\mathbb{R}^2}$ and
the covariance matrix $Q_k$ is

\[Q_k = \begin{pmatrix} a_x^2 & 0 \\
0 & a_y^2
\end{pmatrix}\] 
$F$ is a matrix where

\[F = \begin{pmatrix}
1 & 0 & \Delta t & 0 \\
0 & 1 & 0        & \Delta t \\
0 & 0 & 1        & 0 \\
0 & 0 & 0        & 1
\end{pmatrix}\]
and $G$, also a matrix, is given by

\[G = \begin{pmatrix}
\Delta t^2 & 0 \\
0          & \Delta t^2 \\
\Delta t   & 0 \\
0          & \Delta t
\end{pmatrix}\]

Our observation model becomes

\[\mathbf{z}_k = h_k(\mathbf{x}_k, \mathbf{n}_k) = H\mathbf{x}_k + \mathbf{n}_k\]
where $\mathbf{z}_k = (z_{k_x}, z_{k_y})\tr$, $\mathbf{n}_k$ is a multivariate 
Gaussian variable with mean $\{\mathbf{0}\}_{\mathbb{R}^2}$ and 
covariance $R_k$ (a $2 \times 2$ matrix), and $H$ is a matrix given by

\[H = \begin{pmatrix}
1 & 0 & 0 & 0 \\
0 & 1 & 0 & 0 \\
\end{pmatrix}\]

It is clear that both our state and observation model are linear.

\subsection{Prior and likelihood distributions}

Now that we have these two models in place, we can use them to describe what
in the probability literature are called the ``prior'' distribution and the
``likelihood'' distribution.  A good deal of literature on particle filters 
uses this vocabulary.

The prior distribution is the probability density function that is conditional
on the previous time step.  It is written as 
$p(\mathbf{x}_k|\mathbf{x}_{k-1})$.  This function answers the question: 
given a particular value of $\mathbf{x}_{k-1}$ in the past, what is the 
probability of getting the input $\mathbf{x}_k$ on this time step?  
This can be answered using our function $f_k$.  

In our example, by looking at $f_k$ we can tell that since 
$\mathbf{v}_k$ is Gaussian noise with mean $\{\mathbf{0}\}_{R^{n_v}}$ 
and variance 
$Q_k$, we know that $\mathbf{x}_k \sim \mathcal{N}(F\mathbf{x}_{k-1},GQ_kG\tr)$.
This is because of the property of a multivariate normal distrubtion that if
$Y = \mathbf{c} + BX$, where 
$X \sim \mathcal{N}(\mathbf{\mu}, \Sigma)$, then 
$Y \sim \mathcal{N}(c + B\mathbf{\mu},B\Sigma B\tr)$.  Therefore, we know that

\[p(\mathbf{x}_k|\mathbf{x}_{k-1}) = \frac{1}{(2\pi)^{N_x/2} \det(GQ_kG\tr)^{1/2}} 
e^{-\frac{1}{2}(\mathbf{x}_k - F\mathbf{x}_{k-1})\tr(GQ_kG\tr)^{-1}(\mathbf{x}_k
- F\mathbf{x}_{k-1})} \]

Because this is painful to write out, it is easier to write

\[p(\mathbf{x}_k|\mathbf{x}_{k-1}) = \mathcal{N}(\mathbf{x}_k ; 
F\mathbf{x}_{k-1}, GQ_kG\tr)\]

Here, $\mathbf{x}_k$ is the argument to the pdf of the distribution, 
$F\mathbf{x}_{k-1}$ is the mean, and $GQ_kG\tr$ is the covariance matrix.

The final probability density function we need is the likelihood function.  This
is written as $p(\mathbf{z}_{k} | \mathbf{x}_k)$, and it answers the question:
given that the state is $\mathbf{x}_k$, how likely is the observation
$\mathbf{z}_k$?  We can use $h_k$ to answer this question.

In our example, again we can see from the same argument for $\mathbf{x}_k$
that $\mathbf{z}_k \sim \mathcal{N}(H\mathbf{x_k}, R_k)$.  Thus, we can write

\[p(\mathbf{z}_k|\mathbf{x}_k) = \mathcal{N}(\mathbf{z}_k ; H\mathbf{x}_{k}, R_k)\]

Being able to switch back and forth between the SDE and probability frameworks
will prove very beneficial when reading the literature.  It is worth noting
that in this framework, the problem of determining state from noisy observations
is often stated as trying to find the conditional PDF

\[p(\mathbf{\hat{x}}_k|\mathbf{z}_{1,\ldots,k})\]
where $\mathbf{\hat{x}}_k$ is the estimate of the true state $\mathbf{x}_k$ and 
$\mathbf{z}_{1,\ldots,k}$ are the observations up to time $k$.

\section{Kalman Filter variants}

We now begin by discussing the Kalman filter.  If the state and observation 
models are linear, as in our example, and the noise is Gaussian, the predicted
state $\mathbf{\hat{x}}_k$ given the observations will be Gaussian as well.
The Kalman filter is the optimal predictor of state if the state and 
observation model are linear, and it also optimally estimates the 
covariance of the prediction.  The derivation will not be presented here; it 
can be found in many other places such as Wikipedia and other papers.

The algorithm is presented in the following Matlab code:

\lstinputlisting{kalman.m}

The main point is that the Kalman filter makes a prediction of the state
$m_{k|k-1}$ and covariance $P_{k|k-1}$, then uses this prediction to update
and find the optimal prediction $m_{k|k}$ and the covariance $P_{k|k}$.

\subsection{Extended Kalman Filter}

If the model is not linear, the derivative of the state model is a linear 
approximation of $f_k$ near the predicted point 

\[ F_k = \left.\frac{df_k}{d\mathbf{x}}(\mathbf{x})\right|_{\mathbf{x} = m_{k-1|k-1}} \]

Likewise, the derivative of the observation model is a linear approximation
of $h_k$ near the point $\mathbf{x}_k$.

\[ H_k = \left.\frac{dh_k}{d\mathbf{x}}(\mathbf{x})\right|_{\mathbf{x} = m_{k|k-1}} \]

These approximations can be used with the Kalman filter for a nonlinear 
model, but because the predicted state need not be Gaussian, the extended
Kalman filter is no longer considered optimal in anyway.

\section{Particle filter variants}

Particle filters attempt to estimate the state of a model by running the same
system many times with different paths chosen based on the statistics of the
system.  Each particle has an associated weight, $W$, that tells how likely the 
individual particle is.  The particles must be resampled as they have a
tendancy to become degenerate after a few steps of the model.

The resampling algorithm is shown in the following Matlab code:

\lstinputlisting{resample.m}

The main idea is that we sample a $\theta_k \sim U(0,1)$ and then we pick
$\mathbf{x}_k^* = \mathbf{x}_n$ so that

\[A^{-1} \sum_{j=1}^{n}W_j < \theta_k \le A^{-1}\sum_{j=1}^{n+1} W_j\]
where $A = \sum_{j=1}^{Ns} W_j$.

\subsection{Sampling Importance Resampling Filter}

The code for the sampling importance resampling filter is shown below

\lstinputlisting{particleSir.m}

\subsection{Auxilary Sampling Importance Resampling Filter}

The code for the auxilary sampling importance resampling filter is shown
below

\lstinputlisting{particleAsir.m}

\end{document}
